\documentclass[exercice,a4]{thedocument}%
\startdatakeys
    \source{laclassecooperativedestef.com}
    \BO{2026}
    \cycle{3}
    \competencesDuSocle{RE4,CO2}
    \objectifsApprentissage{C3C1ab,C3C1aj,C3C1ae}
    \motsCles{segment, perpendicularite, quadrilatere}
\enddatakeys
\begin{document}%
\noindent\textbf{Programme de construction 32}\par\smallskip\noindent
\begin{itemize}
    \item Tracer le segment $[\mathspoint{AD}]$ tel que
    $\mathspoint{AD}=10\text{ cm}$.
    \item Placer $\mathspoint{B}$ et $\mathspoint{C}$ sur
    $[\mathspoint{AD}]$ sachant que
    $\mathspoint{AB}=\mathspoint{CD}=2\text{ cm}$ et
    $\mathspoint{BC}=6\text{ cm}$.
    \item Tracer la droite $(\mathspoint{d})$ passant par
    $\mathspoint{B}$, perpendiculaire à $[\mathspoint{AD}]$.
    \item Tracer la droite $(\mathspoint{e})$ passant par
    $\mathspoint{C}$, perpendiculaire à $[\mathspoint{AD}]$.
    \item Tracer $\mathspoint{AEFD}$. Quel est le nom de cette figure ?
\end{itemize}
\end{document}%
